\scanpage{29}

% L1–2
a) $\forall E\in\operatorname{Ob}\Delta$, le foncteur $u\mapsto u\circ\tau_i$ induit une équivalence
\[
 \underline{\mathrm{Hom}}_{d_{i+1}}(C_{i+1},E)-\underline{\mathrm{Hom}}_{d_i}(C_i,E)
\]

% L3
b)
\[
 \tau_i(d_i)\subset d_{i+1}\subset\tau_i(d_i)\cup
 \bigcup_{(S,\lambda)\in\unsure{\rho|\sigma}}\ \bigcup_{u\in\underline{\mathrm{Hom}}_{\deleted{\unsure{\Delta}}}(S,C_{i+1})}u(\lambda)
\]

% L4
c) $\tau_i:C_{\unsure{i|i'}}-C_{i+1}$ est \uline{préexact} pour \deleted{\unsure{les}} $d_i$, i.e.

% L5–6, crossed block and arrow to e)
\deleted{\unsure{d)} $\forall(\unsure{R},R_1,\delta_1,R_2,\delta_2,\varphi\unsure{,\varphi_2})\in\rho$, et $u_1:R_1\to C_i$ qui est $\delta_1$-exact, \inserted{tel que $u_1(\delta_1)\subset d_i$,} il existe un foncteur $\delta_2$-exact $v_2:R_2\to C_{i+1}$ tel que $v_2\circ\varphi_2=(\tau_i u_1)\circ\varphi_1$ \inserted{(et $v_2(\delta_2)\subset d_{i+1}$)}}

% L7–12
d) $\forall(S,\lambda)\in\sigma$ \deleted{\unsure{et}} $u:S\longrightarrow C_i$, \deleted{et et une \unsure{diag|diag.} catégorie $I'$ $(I'')$ intervenant dans $\in\lambda'$ ou $\lambda''$} \inserted{$\unsure{\lambda}$ cône $\in\lambda'$ ou $\lambda''$} \deleted{(i.e. $\unsure{\lambda}=(\lambda',\lambda'')$), il existe un cône \inserted{$c'$} sur le syst. (proj ind) corresp. dans $\tau_i(c)$, tel que $c'\in d_{i+1}$} $\exists$ un foncteur $\bar v:\bar S\longrightarrow C_{i+1}$ (où $\beta:S-\bar S$ est obtenu en dédoublant les sommets des cônes critiques de $S$ $(**)$) tel que $\bar v\beta=\tau_i u$, et tel que $\bar v(\bar\lambda)\subset d_{i+1}$

% L10–12, slanted left
% slanted
\margin{cf définition sur p. $\rightarrow$ avant $(**)$}

% L13
\deleted{Prouvons que pour $\omega$ grand, $(C_\omega,d_\omega)$ est}

% L14–15
e) \quad \begin{math}\left\{\vcenter{\hbox{\begin{minipage}{.87\linewidth}
$e_0)$ $\forall(R_1,\delta_1,-)\in\rho$ et $u_1,v_1:R_1\rightrightarrows C_i$ \inserted{tels que $u_1(\delta_1)\subset d_i$, $v_1(\delta_1)\subset d_i$,} et $\alpha_1,\beta_1:u_1\rightrightarrows v_1$, la relation $\alpha*\varphi=\beta*\varphi$ implique $\tau_i*\alpha=\tau_i*\beta$
% L16–18
$e_1)$ $\forall(R_1,\delta_1,{-}{-})\in\rho$ et $u_1,v_1:R_1\rightrightarrows C_i$ tels que $u_1(\delta_1),v_1(\delta_1)\subset d_i$, et tout hom $u_1\circ\varphi\xrightarrow{\alpha_2}v_1\circ\varphi$, $\exists$ un hom $\tau_i u_1\xrightarrow{\bar\alpha_1}\tau_i v_1$ tel que \deleted{\unsure{e)}} $\bar\alpha_1*\varphi=\tau_i*\alpha_2$
% L19–20
$e_2)$ $\forall(R_1,\delta_1,{-}{-})\in\rho$ et $u_2:R_2-C_i$ tel que $u_{\unsure{2|1}}(\delta_2)\subset d_i$, $\exists\,\bar u_1:R_1\longrightarrow C_{i+1}$ tels que $\bar u_1\circ\varphi=\tau_i\circ u_2$ et $\bar u_1(\delta_1)\subset d_{i+1}$
\end{minipage}}}\right.\end{math}
% archive sheet number: 28
